\chapter{Javascript}
\section{Variabili}
\subsection{var}
Sono variabili che hanno \textbf{function scope}, sono visibili in tutta la funzione che le contiene e \uline{non rispettano lo scope di blocco}.

\begin{lstlisting}
var x = 10; // qui x vale 10
{
  let x = 2; // qui x vale 2
  let y = 5; // qui y vale 5
  var z = 8; // qui z vale 8
}

// qui x vale 10
// y NON è accessibile qui
// z può essere usata qui (var ignora il blocco)
\end{lstlisting}

Si usano per le variabili globali.
\subsection{let}
Sono variabili con \textbf{block scope}, esistono solo nel blocco \texttt{\{\dots\}} in cui è dichiarata.

\subsection{const}
Sono variabili con \textbf{block scope}, esistono solo nel blocco \texttt{\{\dots\}} in cui è dichiarata.

\uline{Non possono essere riassegnate dopo l'inizializzazione.}

Sono contenitori costanti, ma i valori contenuti possono cambiare:
\begin{itemize}
    \item Le variabili semplici sono modificabili.
    \begin{lstlisting}
var x = 10;
// Here x is 10
{
  const x = 2; // Here x is 2
  x = 4;       // ERROR!
}
// Here x is 10
    \end{lstlisting}
    \item I valori delle variabili composte possono cambiare
    \begin{lstlisting}
const ages = [15, 35, 28, 18];
// Here ages is 15, 35, 28, 18

for (let i = 0; i < ages.length; i++) {
  ages[i] = ages[i] + 1;
}
// Here ages is 16,36,29,19
ages[4] = 56;                           
// Here ages is 16,36,29,19,56
ages = [43, 26];
    \end{lstlisting}
\end{itemize}

\section{Controllo di flusso}
\subsection{Strutture standard}
\texttt{if/else, switch, for, while, do while} identici a java
\subsection{Iteratori specifici di JS}
\begin{itemize}
    \item \textbf{for of}: Itera sui valori (array, strignhe, ecc).
    \item \textbf{for in}: Itera sualla chiave di un oggetto.
\end{itemize}


\begin{lstlisting}
const arr = [10, 20, 30];

// for classico
for (let i = 0; i < arr.length; i++) {
  console.log(arr[i]);
}

// itera sui valori
for (const v of arr) {
  console.log(v); // 10, 20, 30
}

// itera sugli indici/chiavi
for (const k in arr) {
  console.log(k); // 0, 1, 2
}
\end{lstlisting}

\newpage

\section{Funzioni}
\subsection{Sintassi di una funzione}
Esistono 3 modi per dichiarare una funzione:
\begin{enumerate}
    \item \textbf{Metodologia classica}
    \begin{lstlisting}
function f(x, y){
    return x * y;
}
    \end{lstlisting}
    \item \textbf{Assegnazione ad una variabile}
    \begin{lstlisting}
var f = function(x,y){
    return x * y;
}
    \end{lstlisting}
    \item \textbf{Arrow function}
    \begin{lstlisting}
const f = (x,y) => {return x*y;}
    \end{lstlisting}

    In questo caso la parola chiave \textbf{\texttt{return}} e le parentesi graffe possono essere omesso solo quando la funzione ha una sola istruazione, per questa ragione è buona abitudine usarle sempre
\end{enumerate}

\section{Stringhe}
\subsection{Definizione}
Una strigna in JS è \uline{una sequenza di 0 o più caratteri racchiusi tra virgolette}
\begin{lstlisting}
let carName1 = "Panda"; // doppie virgolette
let carName2 = 'Punto'; // virgolette singole
\end{lstlisting}

\subsection{Lunghezza di una strigna}
Per ottenere la lunghezza di una strigna, si usa la proprietà length integrata.
\begin{lstlisting}
let sln = carName1.length: //= 5
\end{lstlisting}

\subsection{Concatenazione di stringhe}
L'operatore \texttt{+} può essere usato per concatenare le stringhe.

\begin{lstlisting}
let txt1 = "John";
let txt2 = 'Doe';
let txt3 = txt1 + " " + txt2; // John Doe
\end{lstlisting}

Sommando un numero a una stringa si ottiene una stringa
\begin{lstlisting}
let txt = Ciao;
txt += 35; // Ciao35
\end{lstlisting}

\subsection{Metodi}
\begin{itemize}
    \item Il metodo \textbf{\texttt{indexOf()}} restituisce l'indice (la posizione) della prima occorrenza di un testo specifico in una stringa.

    \begin{lstlisting}
const str = "Please locate where 'locate' occurs!";
let pos = str.indexOf("locate"); // 7
    \end{lstlisting}

    Restituisce -1 se il testo non viene trovato.
    
    Accetta un secondo parametro che indica la posizione di partenza per la ricerca.

    \item Il metodo \textbf{\texttt{lastIndexOf()}} restituisce l'indice dell'\textbf{ultima} occorrenza di un testo specifico in una stringa
    \begin{lstlisting}
let pos = str.lastIndexOf("locate"); // 21
    \end{lstlisting}

    Restituisce -1 se il testo non viene trovato.
    
    Accetta un secondo parametro che indica la posizione di partenza per la ricerca.

    \item Il metodo \textbf{\texttt{startsWith() [endsWith()]}} restituisce \textbf{\texttt{true}} se una strigna inizia [termina] con un valore specificato, altrimenti \textbf{\texttt{false}} (case sensitive).
    \begin{lstlisting}
const str = "Please locate where 'locate' occurs!";
let pos = str.indexOf("locate");
    \end{lstlisting}

    Un secondo parametro opzionale indica una posizione da cui iniziare [a cui terminare]
    \begin{lstlisting}
text.startsWith("world", 6); // controlla dalla posizione 6, restituisce true
text.endsWith("world", 11);  // controlla i primi 11 caratteri, restituisce true
    \end{lstlisting}

    \item Il metodo \textbf{\texttt{slice()}} estrae una porzione di una stringa e restituisce la parte estratta in una nuova stringa.

    Accetta 1 o 2 parametri: la posizione inziiale e la posizione finale (esclusa) [opzionale]
    \begin{lstlisting}
const str = "Apple, Banana, Kiwi";
let pos = str.slice(7, 13); // // Banana
pos = str.slice(7);         // // Banana, Kiwi
    \end{lstlisting}
\end{itemize}

\newpage
\section{Metodi per manipolare Array}
\begin{itemize}
    \item \textbf{Aggiungere elementi}

    Il modo più semplice per aggiungere un nuovo elemento a un array è usare il metodo \textbf{\texttt{push()}}
    \begin{lstlisting}
const fruits = ["Banana","Orange","Apple","Mango"];
fruits.push("Lemon"); // aggiunge un nuovo elemento ("Lemon")
let x = fruits.push(); // restituisce length: il valore di x è 5
    \end{lstlisting}

    \item \textbf{Rimuovere l'ultimo elemento}

    Il metodo \textbf{\texttt{pop()}} rimuove l'ultimo elemento da un array
    \begin{lstlisting}
fruits.pop(); // rimuove l'ultimo elemento ("Lemon")
let y = fruits.pop(); // il valore di y è "Lemon"
    \end{lstlisting}

    \item \textbf{Rimuovere il primo elemento}
    
    Il metodo \texttt{shift()} rimuove il primo elemento dell'array e fa scorrere tutti gli altri elementi a un indice inferiore.
    \begin{lstlisting}
fruits.shift(); // rimuove il primo elemento ("Banana")
let z = fruits.shift(); // il valore di z è "Banana"
    \end{lstlisting}
    \item \textbf{Aggiungere in testa}
    
    Il metodo \texttt{unshift()} rimuove il primo elemento dell'array e fa scorrere tutti gli altri elementi a un indice superiore
    \begin{lstlisting}
fruits.unshift("Banana"); // aggiunge "Banana" in prima posizione
    \end{lstlisting}
    \item \textbf{Splice di un array}

    Il metodo \texttt{splice()} può essere usato per inserire nuovi elementi in un array
    \begin{lstlisting}
const fruits = ["Banana", "Orange", "Apple", "Mango"];
fruits.splice(2, 0, "Lemon", "Kiwi");
// fruits = ["Banana", "Orange", "Lemon", "Kiwi", "Apple", "Mango"]
    \end{lstlisting}
    \begin{itemize}
        \item Il primo parametro indica la posizione in cui i nuovi elementi devono essere aggiunti (inseriti).
        \item Il secondo parametro definisce quanti elementi devono essere rimossi.
        \item I parametri successioni definiscono i nuovi elementi da aggiungere
    \end{itemize}

    Restituisce un array con gli elementi rimossi
    \begin{lstlisting}
let removed = fruits.splice(2, 2, "Lemon", "Kiwi");
// fruits = ["Banana", "Orange", "Lemon", "Kiwi"]
// removed = ["Apple", "Mango"]
    \end{lstlisting}
\end{itemize}

\section{Classi}
\subsection{Definizione}
Permettono di definire dei tipi do oggetti e creare istanze.

In JS tutti i valori sono oggetti, tranne i tipi primitivi.

Quando scriviamo un letterale stiamo creando l'istanza di una classe predefinita
\begin{lstlisting}
// Utilizzo del literal di array (scelta consigliata)
const names = ["Anna", "Diana"];

// Equivalente utilizzando il costruttore Array
const names = new Array("Anna", "Diana");
\end{lstlisting}

\subsection{Anatomia di una classe}
\subsubsection{Elementi di una classe}
\begin{itemize}
    \item \textbf{\texttt{constructor}}: Viene eseguito da new e inizializza i campi tramite this.
    \item \textbf{\texttt{Metodi}}: Si scrivono direttamente nel corpo della classe, senza la parola function.
    \item \textbf{\texttt{Campi}}: Non si dichiararono a priori, nascono assengando \texttt{this.x = \dots} nel constructor
\end{itemize}

\subsubsection{Differenze da Java}
\begin{itemize}
    \item Niente tipi nelle firme
    \item Niente \texttt{public/private} 
    \item Niente overloading: una sola firma per nome
\end{itemize}

\begin{tcolorbox}[title={Esempio di classe}, breakable]
\begin{lstlisting}
class Car {
  // Il costruttore viene chiamato automaticamente quando crei un nuovo oggetto
  constructor(brand, year) {
    this.brand = brand; // Assegna il valore all'istanza dell'oggetto
    this.year = year;
  }

  // Metodo per calcolare l'età dell'auto
  age(currentYear) {
    return currentYear - this.year;
  }

  // Metodo per restituire una stringa descrittiva
  toString() {
    return `${this.brand} (${this.year})`;
  }
}

// Esempio di utilizzo:
// 'new' crea una nuova istanza (oggetto) della classe Car
const c = new Car("Fiat", 2020);

console.log(c.toString()); // Output: "Fiat (2020)"
console.log(c.age(2026));  // Output: 6
\end{lstlisting}    
\end{tcolorbox}

\subsection{Iterazioni su variabili multivalore}
\subsubsection{\texttt{for/of}}
Considera tutti i valori di una struttura iterativa
\begin{lstlisting}
const numbers = [1, 2, 3, 4, 5, 6];
let sum = 0; // Inizializziamo la variabile che conterrà la somma

// Il ciclo for...of scorre ogni elemento dell'array
for (let x of numbers) {
  sum = sum + x; // Aggiunge il valore corrente (x) alla somma totale
}

// Dopo il ciclo, sum vale 21 (1+2+3+4+5+6 = 21)
console.log(sum);
\end{lstlisting}

\subsubsection{\texttt{for/in}}
Considera tutte le proprietà di un oggetto
\begin{lstlisting}
const person = {fname:"John", lname:"Doe", age:25};
let str = "";

// Il ciclo for...in scorre le chiavi dell'oggetto
for (let x in person) {
  // Concatena la chiave (x), il separatore, il valore (person[x]) e un tag HTML
  str += x + ": " + person[x] + "<br>";
}
\end{lstlisting}

\newpage
\section{Programmazione funzionale}
\subsection{\texttt{forEach}}
Il metodo \textbf{\texttt{forEach()}} invoca una funzione (una funzione di callback) una volta per ciascun elemento dell'array 
\begin{lstlisting}
const numbers = [45, 4, 9, 16, 25];
let txt = "";
numbers.forEach(myFunction);

function myFunction(value, index, array) { // index e array sono opzionali
  txt += value + "<br>";
}
\end{lstlisting}

La funzione accetta 3 argomenti:
\begin{itemize}
    \item Il valore dell'elemento
    \item L'indice dell'elemento
    \item L'array stesso
\end{itemize}

\subsection{\texttt{map}}
Il metodo \textbf{\texttt{map()}} crea un nuovo array applicando una funzione a ciascun elemento dell'array.

\begin{itemize}
    \item Non esegue la funzione sugli elementi privi di valore.
    \item Non modifica l'array originale
\end{itemize}

\begin{lstlisting}
const numbers1 = [45, 4, 9, 16, 25];
const numbers2 = numbers1.map(myFunction);

function myFunction(value, index, array) { // moltiplica ogni valore per 2
  return value * 2;
} // array risultante: [90, 8, 18, 32, 50]
\end{lstlisting}

La funzione accetta 3 argomenti:
\begin{itemize}
    \item Il valore dell'elemento.
    \item L'indice dell'elemento.
    \item L'array stesso.
\end{itemize}

\subsection{\texttt{filter}}
Il metodo \textbf{\texttt{filter()}} crea un nuovo array contente gli elementi che superano un test
\begin{itemize}
    \item Non esegue la funzione sugli elementi senza valore.
    \item Non modifica l'array originale.
\end{itemize}
\begin{lstlisting}
const numbers = [45, 4, 9, 16, 25];
const over18 = numbers.filter(myFunction);

function myFunction(value, index, array) {
  // crea un nuovo array con gli elementi maggiori di 18
  return value > 18;
} // array risultante: [45, 25]
\end{lstlisting}


La funzione accetta 3 argomenti:
\begin{itemize}
    \item Il valore dell'elemento.
    \item L'indice dell'elemento.
    \item L'array stesso.
\end{itemize}

\subsection{\texttt{reduce}}
Il metodo \textbf{\texttt{reduce()}} esegue una funzione su ogni elemento dell'array per produrre un singolo valore.
\begin{itemize}
    \item Lavora da sinistra a destra sull'array
    \item Non modifica l'array originale
\end{itemize}
\begin{lstlisting}
const numbers = [45, 4, 9, 16, 25];
let sum = numbers.reduce(myFunction);

function myFunction(total, value, index, array) { // somma tutti i numeri di un array
  return total + value;
} // valore risultante: 99
\end{lstlisting}

La funzione accetta 4 argomenti:
\begin{itemize}
    \item L'accumulatore (il valore iniziale o il valore restituito in precedenza)
    \item Il valore corrente
    \item L'indice corrente
    \item L'array stesso
\end{itemize}

\section{Visualizzare i dati}
JS non ha funzioni native di stampa, per "mostrare" qualcosa si scrive nel DOM, in una finiestra di dialogo o nella console

\subsection{Scrivere nel DOM (modo principale)}
\begin{itemize}
    \item \texttt{element.textContent}: Inserisce solo testo, sicuro contro XSS
    \item \texttt{element.innerHTML}: Interpreta markup HTML, usare solo con contenuto fidato
\end{itemize}
\subsection{Console del browser (debug)}
\begin{itemize}
    \item \texttt{console.log()}: stampa nella DevTools console
\end{itemize}
\subsection{Finiestre di dialogo (uso limitato)}
\begin{itemize}
    \item \texttt{alert(), confirm(), prompt()}: Bloccano l'interfaccia, usare solo per i prototipi o casi eccezionali  
\end{itemize}

\section{Gestione eventi}
Javascript collega una funzione (handler) all'evento, e quel codice viene eseguito quando l'evento si verifica

\subsection{Cos'è un evento?}
Click, input da tastiera, submit di un form, caricamento pagina, scroll, mouseover, ...

\subsection{Come aggiungere un Listener}
Si usa il metodo \texttt{element.addEventListener(event, handler, options)}
\begin{itemize}
    \item \textbf{tipo di evento}
    \begin{itemize}
        \item Stringa: "click", "mousedown", "keyup", ...
        \item Senza prefisso "on": "click", non "onclick"
    \end{itemize}
    \item \textbf{handler}
    \begin{itemize}
        \item Funzione anonima oppure nome di funzione (senza parametri)
        \item Riceve l'oggetto Event come argomento
    \end{itemize}
    \item \textbf{(opzionale) - options}
    \begin{itemize}
        \item \{once: true\}, \{capture: true\}
    \end{itemize}
\end{itemize}

\begin{tcolorbox}[title={Esempio}]
\begin{lstlisting}[language=html]
<!DOCTYPE html>
<html>
<body>
    <button id="saveBtn">Salva</button>
    <script>
        const btn = document.getElementById("saveBtn");

        // Handler con funzione anonima
        btn.addEventListener("click", function (e) {
            console.log("Cliccato!", e.target);
        });

        // Handler nominato + options
        function onSave(e) { /* .. */ }
        btn.addEventListener("click", onSave, { once: true });
    </script>
</body>
</html>
\end{lstlisting}
\end{tcolorbox}
